۳۰ ۳۱

شاید به این دلیل بود که جذب جنس مخالف تنها زمینه‌ای از زندگی‌ام بود که در آن احساس شکست کامل داشتم. هر بار که در خیابان قدم می‌زدم یا وارد بار می‌شدم، شکست خودم را دیدم که با رژ لب قرمز و ریمل مشکی به من خیره شده است. ترکیب اشتیاق و فلج شدن مرگبار بود.

بعد از کارگاه آن شب، کابینت فایل خود را باز کردم و در میان کاغذهایم جستجو کردم. چیزی بود که می‌خواستم پیدا کنم، چیزی که سال‌هاست به آن نگاه نکرده بودم. بعد از نیم‌ساعت، آن را پیدا کردم: پوشه‌ای با برچسب «نوشته‌های دبیرستان». برگه‌ای از دفتر خط‌دار را بیرون کشیدم که از بالا تا پایین پر از خرچنگ‌نویسی‌هایم بود. تنها شعری بود که در زندگی‌ام سعی کرده‌ام. در کلاس یازدهم نوشته شده بود و هرگز به کسی نشان ندادم. با این حال، جواب سوالم بود.

ناامیدی جنسی

اثر نیل استراوس

تنها دلیلی که به جایی می‌روی، برای شلیک به هدفی که به ندرت به آن می‌رسی. تنها هدف در ذهن، نگاه گذرا به یک جفت پای آشنا در خیابانی شلوغ یا فشاری از یک زن که فقط می‌توانی او را دوست بنامی.

شب بی‌امتیاز خصومت ایجاد می‌کند. آخر هفته بی‌امتیاز بدبینی به بار می‌آورد. حتی مرد خردمند در بهشت احمقان زندگی می‌کند. از چشمانی قرمز تمام جهان دیده می‌شود، عصبانیت از دوستان و خانواده به دلایل نامعلومی برایشان. فقط خودت می‌دانی چرا آنقدر عصبانی هستی.

آن دیگری که صرفاً دوستت است، که برای مدت طولانی می‌شناختی، که تو را خیلی احترام می‌کند که نمی‌توانی آنچه را که می‌خواهی انجام دهی. و او دیگر زحمت نمی‌دهد خودش را نقاب‌بزنانه نشان دهد و معاشقه کند چون فکر می‌کند که او را همان‌طور که هست دوست داری، در حالی که چیزی که در مورد او دوست داشتی، معاشقه‌اش بود.

در دهه‌ای که از نوشتن آن شعر گذشته بود، هیچ چیز تغییر نکرده بود. هنوز هم نمی‌توانستم شعر بنویسم. و مهم‌تر از آن، هنوز هم احساس مشابهی داشتم. شاید ثبت‌نام در کارگاه مستری تصمیمی هوشمندانه بوده باشد. به هر حال، داشتم کاری فعالانه برای بهبود ناتوانی‌ام انجام می‌دادم.